\section{EXPERIMENTAL METHOD}
In this experiment, students must carry out the following tasks:
\begin{enumerate}
    \item Determine the status of air that includes: temperature, humidity behind the cooler, before the air drying, before the steam nozzle and after the steam nozzle. From the measured data, students must draw changing status of air on I-d diagram.
    \item Calculate the heat balance of the aerodynamic tube including: determining the flowrate of air, determining the cold yield of the cooler and the heating load of the drying equipment.
\end{enumerate}

\subsection{Operation process}
\begin{enumerate}
    \item Turn on the gate switch, check the 3-phase indicator on the main switchboard.
    \item Turn on fan, adjust air flowrate by opening/closing the door.
    \item Turn on air conditioner.
    \item Turn on resistor.
    \item Turn on the steam control for saturated steam. Observe temperature and pressure in the steam tank, start the steam valve when the pressure reaches 2 kg/cm$^2$.
    \item Let the system run for 2 minutes to achieve stability. Measure dry and wet bulb temperatures at locations. Measure the condensed water flow.
    \item Turn on the steam control for superheated steam, let the system run for 2 minutes and do measurement like above.
    \item Change the other operating mode by changing the wind outlet position, increasing or decreasing the resistor, or changing the airflow.
\end{enumerate}

\vspace{0.3cm}
\noindent \textbf{Note:}
The water level in the boiler was tested after each experiment (turn off the resistor) by opening/closing the valve between the boiler and the water tank to provide water for the boiler. Provided water level is equal to saturated steam thermometer.
